\documentclass{article}
\usepackage{amsmath,amssymb,amsthm}

\newtheorem{theorem}{Theorem}

\begin{document}

\begin{theorem}
For every natural number $n \in \mathbb{N}$, there exists a prime number
$p$ such that $n \le p$. Consequently, there are infinitely many prime
numbers.
\end{theorem}

\begin{proof}
Let $n \in \mathbb{N}$ be given. Consider the natural number $n! + 1$.
Since $n! \ge 1$, we have $n! + 1 \ge 2$, so $n! + 1 \neq 1$.

Let $p$ be the smallest prime factor of $n! + 1$. Since $n! + 1 \neq 1$,
such a prime factor exists, and by definition $p$ is prime.

We now claim that $n \le p$. Suppose, for the sake of contradiction, that
$p < n$.
Since $p$ is a positive prime with $p \le n$, $p$ divides the factorial $n!$
(that is, $p \mid n!$). Furthermore, by definition of $p$, we have
$p \mid (n! + 1)$.

Since $p$ divides both $n!$ and $n! + 1$, $p$ must also divide their
difference:
\[
p \mid \bigl((n! + 1) - n!\bigr) \implies p \mid 1.
\]
However, no prime number divides $1$, yielding a contradiction.

Hence, our assumption that $p < n$ must be false, so $n \le p$. Thus, $p$
is a prime number satisfying $n \le p$, as required.
\end{proof}

\end{document}
